% GENERATED FILE. Edit canonical lessons, then run npm run generate:books.
% canonical-source-hash: fnv1a64:32b88a46f12789d5
% canonical-lessons: FR-C453-emballage, FR-C453-emballer, FR-C453-livreur, FR-C453-carton, FR-C453-recommande

\chapter{Emballer et envoyer}
\label{ch:fr-emballer-et-envoyer}
\hlchaptermodality{453}

\begin{chapteropening}
\textbf{By the end of this chapter:} \emph{I can send things: the packaging, wrapping, the delivery man, the cardboard box and registered post.}

Complete the last lesson of chapter 453: I can send things: the packaging, wrapping, the delivery man, the cardboard box and registered post.
\end{chapteropening}

\section[l'emballage]{l'emballage — the packaging, the wrapping}

\label{lesson:FR-C453-emballage}



\begin{quote}
Before the new one: say the French for to sell off cheap, then the French for the street sale.
\end{quote}

\subsection*{You'll want to know: l'emballage}
\textbf{l'emballage} — \textquotedblleft{}the packaging, the wrapping\textquotedblright{}.

\begin{quote}
Il y a trop d'emballage autour de ce petit objet.
\end{quote}

\begin{quote}
\emph{There is too much packaging round this small object.}
\end{quote}

A masculine noun. \emph{Un emballage cadeau} — gift wrap.

\begin{cousinweb}
\emph{emballage} is built on \emph{emballer} 'to wrap up' with \emph{-age} (EtymDB lists it). \emph{Emballer} comes from \emph{balle} 'bale', usually traced to a Frankish word.
\end{cousinweb}

\begin{grammarlens}[title={What you've built}]
One more word for telling, explaining and describing at B1.
\end{grammarlens}

\subsection*{Your turn}
\begin{itemize}
\raggedright
  \item \emph{Say it:} \emph{l'emballage}
  \item \emph{Say it:} \emph{l'emballage}, once more
  \item \emph{Say it:} a sentence of your own with \emph{l'emballage}
  \item \emph{Recall:} say the French for to sell off cheap, then the French for the street sale, then say \emph{l'emballage} again
\end{itemize}

\begin{tcolorbox}[breakable,colback=teal!4,colframe=teal!35!black,title={Before you move on}]
What does l'emballage mean? (\textquotedblleft{}The packaging, the wrapping\textquotedblright{}.) Say it once more.
\end{tcolorbox}
\section[emballer]{emballer — to wrap (up), to pack; to thrill (informal)}

\label{lesson:FR-C453-emballer}



\begin{quote}
Before the new one: say the French for the street sale, then the French for the packaging.
\end{quote}

\subsection*{You'll want to know: emballer}
\textbf{emballer} — \textquotedblleft{}to wrap (up), to pack; to thrill (informal)\textquotedblright{}.

\begin{quote}
La vendeuse emballe le cadeau dans du papier rouge.
\end{quote}

\begin{quote}
\emph{The shop assistant wraps the present in red paper.}
\end{quote}

Informal: \emph{ça m'emballe} — I love the idea; \emph{s'emballer} — to get carried away.

\begin{cousinweb}
\emph{emballer} is built from \emph{en-} and \emph{balle} 'bale, bundle' (EtymDB lists the verb), and \emph{balle} is usually traced to a Frankish word meaning 'ball'.
\end{cousinweb}

\begin{grammarlens}[title={What you've built}]
One more word for telling, explaining and describing at B1.
\end{grammarlens}

\subsection*{Your turn}
\begin{itemize}
\raggedright
  \item \emph{Say it:} \emph{emballer}
  \item \emph{Say it:} \emph{emballer}, once more
  \item \emph{Say it:} a sentence of your own with \emph{emballer}
  \item \emph{Recall:} say the French for the street sale, then the French for the packaging, then say \emph{emballer} again
\end{itemize}

\begin{tcolorbox}[breakable,colback=teal!4,colframe=teal!35!black,title={Before you move on}]
What does emballer mean? (\textquotedblleft{}To wrap (up), to pack; to thrill (informal)\textquotedblright{}.) Say it once more.
\end{tcolorbox}
\section[le livreur]{le livreur — the delivery man, the courier}

\label{lesson:FR-C453-livreur}



\begin{quote}
Before the new one: say the French for the packaging, then the French for to wrap.
\end{quote}

\subsection*{You'll want to know: le livreur}
\textbf{le livreur} — \textquotedblleft{}the delivery man, the courier\textquotedblright{}.

\begin{quote}
Le livreur a laissé le paquet chez la voisine.
\end{quote}

\begin{quote}
\emph{The delivery man left the parcel with the neighbour.}
\end{quote}

Feminine \emph{la livreuse}.

\begin{cousinweb}
\emph{livreur} is built on \emph{livrer} 'to deliver' with \emph{-eur} (EtymDB); \emph{livrer} is inherited from Latin \emph{līberāre} 'to set free', from \emph{līber} 'free'.
\end{cousinweb}

\begin{grammarlens}[title={What you've built}]
One more word for telling, explaining and describing at B1.
\end{grammarlens}

\subsection*{Your turn}
\begin{itemize}
\raggedright
  \item \emph{Say it:} \emph{le livreur}
  \item \emph{Say it:} \emph{le livreur}, once more
  \item \emph{Say it:} a sentence of your own with \emph{le livreur}
  \item \emph{Recall:} say the French for the packaging, then the French for to wrap, then say \emph{le livreur} again
\end{itemize}

\begin{tcolorbox}[breakable,colback=teal!4,colframe=teal!35!black,title={Before you move on}]
What does le livreur mean? (\textquotedblleft{}The delivery man, the courier\textquotedblright{}.) Say it once more.
\end{tcolorbox}
\section[le carton]{le carton — the cardboard box; cardboard}

\label{lesson:FR-C453-carton}



\begin{quote}
Before the new one: say the French for to wrap, then the French for the delivery man.
\end{quote}

\subsection*{You'll want to know: le carton}
\textbf{le carton} — \textquotedblleft{}the cardboard box; cardboard\textquotedblright{}.

\begin{quote}
Nous avons mis les livres dans des cartons.
\end{quote}

\begin{quote}
\emph{We put the books in cardboard boxes.}
\end{quote}

\emph{Un carton jaune} — a yellow card, in sport.

\begin{cousinweb}
\emph{carton} was borrowed from Italian \emph{cartone} 'thick paper' (EtymDB), from \emph{carta}, which goes back to Latin \emph{charta} 'sheet of papyrus'.
\end{cousinweb}

\begin{grammarlens}[title={What you've built}]
One more word for telling, explaining and describing at B1.
\end{grammarlens}

\subsection*{Your turn}
\begin{itemize}
\raggedright
  \item \emph{Say it:} \emph{le carton}
  \item \emph{Say it:} \emph{le carton}, once more
  \item \emph{Say it:} a sentence of your own with \emph{le carton}
  \item \emph{Recall:} say the French for to wrap, then the French for the delivery man, then say \emph{le carton} again
\end{itemize}

\begin{tcolorbox}[breakable,colback=teal!4,colframe=teal!35!black,title={Before you move on}]
What does le carton mean? (\textquotedblleft{}The cardboard box; cardboard\textquotedblright{}.) Say it once more.
\end{tcolorbox}
\section[le recommandé]{le recommandé — the registered letter, registered post}

\label{lesson:FR-C453-recommande}



\begin{quote}
Before the new one: say the French for the delivery man, then the French for the cardboard box.
\end{quote}

\subsection*{You'll want to know: le recommandé}
\textbf{le recommandé} — \textquotedblleft{}the registered letter, registered post\textquotedblright{}.

\begin{quote}
La banque m'a envoyé un recommandé.
\end{quote}

\begin{quote}
\emph{The bank sent me a registered letter.}
\end{quote}

\emph{En recommandé} — by registered post; \emph{une lettre recommandée}.

\begin{cousinweb}
\emph{recommandé} is the past participle of \emph{recommander} 'to entrust', from \emph{re-} and \emph{commander}, which goes back to Latin \emph{mandāre} 'to entrust' (Lewis \& Short). EtymDB does not list the noun.
\end{cousinweb}

\begin{grammarlens}[title={What you've built}]
One more word for telling, explaining and describing at B1.
\end{grammarlens}

\subsection*{Your turn}
\begin{itemize}
\raggedright
  \item \emph{Say it:} \emph{le recommandé}
  \item \emph{Say it:} \emph{le recommandé}, once more
  \item \emph{Say it:} a sentence of your own with \emph{le recommandé}
  \item \emph{Recall:} say the French for the delivery man, then the French for the cardboard box, then say \emph{le recommandé} again
\end{itemize}

\begin{tcolorbox}[breakable,colback=teal!4,colframe=teal!35!black,title={Before you move on}]
What does le recommandé mean? (\textquotedblleft{}The registered letter, registered post\textquotedblright{}.) Say it once more.
\end{tcolorbox}
